% =====================================================================
%  Guida allo studio — come usare gli appunti, in due pagine grafiche
%  Compilazione:  powershell -File compila.ps1 guida-studio
% =====================================================================
\documentclass[11pt]{article}
\usepackage[a4paper, landscape, hmargin=1.1cm, vmargin=0.8cm]{geometry}
\usepackage[T1]{fontenc}
\usepackage[utf8]{inputenc}
\usepackage[italian]{babel}
\usepackage{newtxtext}
\usepackage{newtxmath}
\usepackage[scaled=0.88]{beramono}
\usepackage[dvipsnames]{xcolor}
\definecolor{unical}{HTML}{B00000}
\definecolor{blu}{HTML}{0B3A91}
\definecolor{verde}{HTML}{1E7B34}
\definecolor{arancio}{HTML}{C0620A}
\definecolor{grigio}{HTML}{4A4A4A}
\definecolor{sfondo}{HTML}{F5F6F8}
\usepackage{tikz}
\usetikzlibrary{arrows.meta, positioning, calc}
\usepackage[hidelinks]{hyperref}
\hypersetup{pdftitle={Algorithmic Game Theory -- Guida allo studio}, pdflang={it}}
\pagestyle{empty}
\setlength{\parindent}{0pt}

\tikzset{
  titolo/.style={font=\sffamily\bfseries\LARGE, text=blu, anchor=north west},
  sottotitolo/.style={font=\sffamily\small, text=grigio, anchor=north west},
  sezione/.style={font=\sffamily\bfseries\large, text=blu, anchor=south west},
  passo/.style={rounded corners=3pt, fill=sfondo, draw=blu!35, line width=0.6pt,
    text width=11.9cm, minimum height=2.45cm, inner sep=7pt, align=left,
    font=\normalsize, anchor=west},
  numero/.style={circle, fill=blu, text=white, font=\sffamily\bfseries\Large,
    minimum size=1.05cm, inner sep=0pt},
  etichetta/.style={font=\sffamily\bfseries\scriptsize, text=white, inner xsep=4pt,
    inner ysep=2.5pt, rounded corners=1.5pt},
  scheda/.style={rounded corners=3pt, draw=grigio!40, fill=white, text width=3.25cm,
    minimum height=2.2cm, inner sep=6pt, align=left, font=\small, anchor=north west},
  freccia/.style={-{Stealth[length=7pt]}, line width=1.1pt, color=blu!70},
}
\newcommand{\file}[1]{\texttt{\small #1}}

% Riga della legenda dei box: y, colore bordo, sfondo, nome, descrizione, priorità, colore priorità
\newcommand{\legenda}[7]{%
  \fill[#3] (16.4,#1) rectangle ++(11.0,-1.12);
  \fill[#2] (16.4,#1) rectangle ++(0.11,-1.12);
  \node[anchor=north west, font=\sffamily\bfseries\small, text=#2] at (16.65,#1-0.1) {#4};
  \node[anchor=north west, text width=6.6cm, font=\small, align=left] at (16.65,#1-0.48) {#5};
  \node[etichetta, fill=#7, anchor=east] at (27.25,#1-0.56) {#6};
}

\begin{document}
% =====================================================================
%  PAGINA 1 — il metodo
% =====================================================================
\noindent
\begin{tikzpicture}[x=1cm, y=1cm]
  \useasboundingbox (0,0) rectangle (27.5,18.8);

  \node[titolo] at (0,18.8) {Come studiare dagli appunti};
  \node[sottotitolo] at (0.02,17.85) {Algorithmic Game Theory --- UniCal, a.a.\ 2026/2027.
    Il metodo per ogni lezione, i box degli appunti e quale PDF aprire.};

  % ---------------- i cinque passi ----------------
  \node[sezione] at (0,16.55) {Per ogni lezione, cinque passi};

  \foreach \n/\y/\tit/\testo in {
    1/15.2/Orientati/{Leggi prima il box \textbf{Riepilogo} in fondo alla lezione: in cinque
      minuti sai dove si va. Poi apri accanto le slide del docente, in \file{lessons/}.},
    2/12.45/Teoria/{Leggi gli appunti nell'ordine. \textbf{Definizioni} e \textbf{Teoremi}
      sono quello che devi saper ripetere; gli \textbf{Esempi} rifalli con carta e penna
      \emph{prima} di leggere la soluzione.},
    3/9.7/Punti critici/{Rileggi i box arancioni \textbf{Attenzione}: segnalano errori nelle
      slide e nei notebook, e i passaggi dove è facile sbagliare. È lì che si vede se hai
      capito davvero.},
    4/6.95/Laboratorio/{Apri il notebook vuoto in \file{labs/NN-.../} e prova a completarlo
      da solo. Poi confronta con il codice degli appunti e con i \textbf{risultati misurati}:
      i tuoi numeri devono tornare.},
    5/4.2/Progetto/{Apri \file{project.pdf}: ogni lezione ``sblocca'' nuove idee nella
      \textbf{banca delle idee}. Annota quale ti convince e quali domande fare al docente.}
  }{
    \node[numero] (n\n) at (0.6,\y) {\n};
    \node[passo] (p\n) at (1.45,\y) {{\sffamily\bfseries\color{blu}\tit}\\[2pt]\testo};
  }
  \foreach \a/\b in {1/2, 2/3, 3/4, 4/5} \draw[freccia] (n\a.south) -- (n\b.north);

  % ritorno alla lezione successiva
  \draw[freccia, dashed, color=grigio!70] (p5.east) -- ++(0.9,0) |- (p1.east)
    node[pos=0.25, right, font=\sffamily\scriptsize, text=grigio, align=left]
    {lezione\\successiva};

  % ---------------- legenda dei box ----------------
  \node[sezione] at (16.4,16.55) {I box degli appunti};
  \legenda{16.3}{blu}{blu!5}{Definizione}{L'enunciato preciso di un concetto.}{ESSENZIALE}{unical}
  \legenda{15.03}{unical}{unical!5}{Teorema}{Enunciato e idea della dimostrazione.}{ESSENZIALE}{unical}
  \legenda{13.76}{verde}{verde!5}{Esempio}{Rifallo a mano: è l'allenamento per l'esame.}{IMPORTANTE}{arancio}
  \legenda{12.49}{arancio}{arancio!7}{Attenzione}{Errori di slide e notebook, trappole tipiche.}{IMPORTANTE}{arancio}
  \legenda{11.22}{grigio}{sfondo}{Osservazione}{Perché le cose funzionano come funzionano.}{IMPORTANTE}{arancio}
  \legenda{9.95}{grigio!60}{white}{Approfondimento}{Oltre le slide: da leggere se hai tempo.}{FACOLTATIVO}{grigio!70}
  \legenda{8.68}{blu}{blu!4}{Riepilogo}{L'intera lezione in poche righe.}{RIPASSO}{blu}

  \node[anchor=north west, font=\small, text width=10.8cm, align=left] at (16.4,7.35)
    {Ogni lezione ha due metà:
     \tikz[baseline=(t.base)]\node[etichetta, fill=blu](t){TEORIA}; e
     \tikz[baseline=(t.base)]\node[etichetta, fill=unical](t){LAB}; (codice verificato,
     con i numeri misurati).};

  % ---------------- quale PDF ----------------
  \node[sezione] at (16.4,5.55) {Quale PDF aprire};
  \node[scheda] (s1) at (16.4,5.45)
    {{\sffamily\bfseries\color{blu}lezione-NN-\ldots.pdf}\\[2pt]
     Per studiare una lezione alla volta, con le slide accanto.};
  \node[scheda, right=0.18cm of s1] (s2)
    {{\sffamily\bfseries\color{blu}main.pdf}\\[2pt]
     Tutto il corso con un indice unico: per il ripasso.};
  \node[scheda, right=0.18cm of s2] (s3)
    {{\sffamily\bfseries\color{unical}project.pdf}\\[2pt]
     Metodo, idee, domande e decisioni per il progetto d'esame.};
  \node[anchor=north west, font=\sffamily\scriptsize, text=grigio] at (16.4,3.05)
    {Tutti in \file{documentazione/pdf/}.};

  % ---------------- se hai poco tempo ----------------
  \node[rounded corners=3pt, fill=unical!6, draw=unical!50, text width=26.8cm,
        inner sep=7pt, anchor=south west, align=left, font=\small] at (0,0.1)
    {{\sffamily\bfseries\color{unical}Se hai poco tempo.}\quad
     Leggi i \textbf{Riepiloghi} di tutte le lezioni, poi i box \textbf{Attenzione}, poi
     rifai i laboratori dal 3 al 6: il loro codice (best response, minimax e equilibri correlati
     con programmazione lineare, core, Shapley esatto e Monte Carlo) è quello che più
     probabilmente riuserai nel progetto. L'esame è un progetto: saper \emph{modellare} e \emph{implementare} conta più
     che ricordare le dimostrazioni.};
\end{tikzpicture}

\newpage
% =====================================================================
%  PAGINA 2 — la mappa del corso
% =====================================================================
\noindent
\begin{tikzpicture}[x=1cm, y=1cm,
  testa/.style={fill=blu, text=white, text width=4.0cm, inner sep=5pt, align=left,
    font=\sffamily\small, minimum height=1.05cm, anchor=north},
  corpo/.style={fill=blu!5, text width=4.0cm, inner sep=5pt, align=left,
    font=\small, minimum height=3.55cm, anchor=north},
  lab/.style={fill=unical!6, draw=unical!45, text width=4.0cm, inner sep=5pt,
    align=left, font=\small, minimum height=1.25cm, anchor=north},
  dipendenza/.style={-{Stealth[length=6pt]}, line width=0.9pt, color=blu!60},
  punto/.style={circle, fill=#1, minimum size=0.32cm, inner sep=0pt},
  nota/.style={font=\scriptsize, text width=4.2cm, align=center, anchor=north},
]
  \useasboundingbox (0,0) rectangle (27.5,18.8);

  \node[titolo] at (0,18.8) {La mappa del corso};
  \node[sottotitolo] at (0.02,17.85) {Che cosa c'è in ogni lezione, che cosa usa delle
    precedenti, e i fili che le attraversano.};

  % ---------------- colonne delle lezioni ----------------
  \foreach \i/\x/\tit/\usa/\punti/\lab in {
    1/2.3/Introduzione all'AGT/{punto di partenza}/
      {Le aree dell'AGT\\Dilemma del prigioniero\\Giochi combinatori\\Minimax e alpha-beta}/
      {Tris: Minimax e alpha-beta},
    2/6.88/Teoria dell'utilità/{indipendente dalla L1}/
      {Preferenze, utilità ordinale\\Lotterie e teorema VNM\\Allais, avversione al rischio\\Ricerca bayesiana}/
      {Battaglia navale: ricerca bayesiana},
    3/11.46/Giochi non cooperativi/{L1 (minimax), L2 (utilità attesa)}/
      {Dominanza e IESDS\\Nash puro e misto\\Condizione del supporto\\Minimax nei giochi a somma zero}/
      {Nash, minimax con PL, label propagation},
    4/16.04/Equilibri e learning/{L3 (Nash, strategie miste, PL)}/
      {Esistenza: Brouwer\\Complessità: PPAD, NP\\Equilibri correlati\\Fictitious play, regret matching}/
      {Equilibri correlati, apprendimento},
    5/20.62/Giochi coalizionali/{L1 (Shapley), L4 (cooperazione)}/
      {Struttura di coalizione\\Core e superadditività\\Nucleolo\\Valore di Shapley}/
      {Core, Shapley, nucleolo},
    6/25.2/Coalizionali: il calcolo/{L5 (Shapley e i suoi assiomi)}/
      {Gioco dell'aeroporto\\Sottografo indotto (ISG)\\MC-Nets\\Shapley Monte Carlo}/
      {Aeroporto, Monte Carlo, ISG, MC-Nets}
  }{
    \node[testa] (t\i) at (\x,14.85)
      {{\scriptsize LEZIONE \i}\\[-1pt]\textbf{\tit}};
    \node[corpo] (c\i) at (t\i.south)
      {{\scriptsize\color{grigio}\textit{Usa:} \usa}\\[4pt]\punti};
    \node[lab] (l\i) at ($(c\i.south)+(0,-0.22)$)
      {{\sffamily\bfseries\color{unical}LAB \i}\quad \lab};
  }

  % ---------------- dipendenze ----------------
  \foreach \a/\b in {2/3, 3/4, 4/5, 5/6}
    \draw[dipendenza] ($(t\a.north east)+(-0.8,0)$) to[bend left=60]
      ($(t\b.north west)+(0.8,0)$);
  \draw[dipendenza] (t1.north) |- (11.46,15.95)
    node[pos=0.75, above, font=\sffamily\scriptsize, text=blu]
    {minimax: dalle mosse a turni alle mosse simultanee} -- (t3.north);
  \draw[dipendenza, dashed] ($(t1.north west)+(0.4,0)$) |- (20.62,16.85)
    node[pos=0.75, above, font=\sffamily\scriptsize, text=blu]
    {Shapley: anticipato nella lezione 1, formalizzato nella 5} -- (t5.north);

  % ---------------- fili conduttori ----------------
  \node[sezione] at (0,7.75) {Fili conduttori};

  % \filo{y}{colore}{nome}: la linea; \tappa{y}{colore}{x}{nota}: una lezione sul filo
  \newcommand{\filo}[3]{%
    \node[anchor=south west, font=\sffamily\bfseries\small, text=#2] at (0,#1+0.08) {#3};
    \draw[line width=1.6pt, color=#2!35] (0.1,#1) -- (27.4,#1);}
  \newcommand{\tappa}[4]{%
    \node[punto=#2] at (#3,#1) {};
    \node[nota] at (#3,#1-0.2) {#4};}

  \filo{6.65}{verde}{Probabilità e incertezza}
  \tappa{6.65}{verde}{6.88}{lotterie, Bayes}
  \tappa{6.65}{verde}{11.46}{strategie miste}
  \tappa{6.65}{verde}{16.04}{distribuzioni sui profili}
  \tappa{6.65}{verde}{25.2}{permutazioni campionate}

  \filo{4.95}{arancio}{Programmazione lineare}
  \tappa{4.95}{arancio}{11.46}{minimax}
  \tappa{4.95}{arancio}{16.04}{equilibri correlati}
  \tappa{4.95}{arancio}{20.62}{nucleolo, core vuoto}

  \filo{3.25}{unical}{Complessità di calcolo}
  \tappa{3.25}{unical}{2.3}{alpha-beta: 549\,946 $\to$ 18\,297 nodi}
  \tappa{3.25}{unical}{16.04}{PPAD, NP}
  \tappa{3.25}{unical}{20.62}{$B_n$, $O(3^n)$, Shapley $2^n$}
  \tappa{3.25}{unical}{25.2}{Shapley in $O(n)$; errore $1/\sqrt{m}$}

  % ---------------- in arrivo ----------------
  \node[rounded corners=3pt, draw=grigio!50, dashed, fill=white, text width=26.8cm,
        inner sep=7pt, anchor=south west, align=left, font=\small] at (0,0.1)
    {{\sffamily\bfseries\color{grigio}In arrivo nel corso.}\quad
     Scelta sociale computazionale (voto, Condorcet)\ \textbullet\ Matching markets
     (matching stabile)\ \textbullet\ Mechanism design (aste)\ \textbullet\ Fair division.
     Gli appunti e la mappa si aggiornano a ogni lezione nuova.};
\end{tikzpicture}

\end{document}
